% Source: source/普通高考/2012/2012福建理.pdf, page 1, question 12.
% Complete node -> directed-edge list from the original PDF:
%   start -> k=1,s=1 -> k<4; k<4 --是--> s=2s-k -> k=k+1
%   -> return to the main stem above the test; k<4 --否--> 输出s -> 结束.
% The return arrow points left into the stem, while the stem's only downward
% arrow enters the test node.
% Grid evidence: tmp/review_2012_fujian_science/grid/page1_grid100.png.
% No-grid source crop: tmp/review_2012_fujian_science/crops/q12_source_tight.png.
% Solid segments: all node borders/connectors; dashed segments: none.
% Branch labels follow the source placement (是 above the right branch, 否 on
% the right of the downward branch) and use direct edge anchors; node text is
% centered with local zero indentation.

\begin{tikzpicture}[
  >=Stealth,
  x=1cm,
  y=0.82cm,
  line width=0.45pt,
  line join=round,
  line cap=round,
  font=\normalsize,
  flowtext/.style={anchor=center, align=center,
    execute at begin node={\setlength{\parindent}{0pt}}},
  every node/.style={flowtext},
  flowbox/.style={flowtext, draw, minimum height=0.68cm,
    inner xsep=4pt, inner ysep=2.5pt},
  branchlabel/.style={flowtext, inner sep=0pt},
  terminal/.style={flowbox, rounded corners=0.18cm, minimum width=1.35cm},
  process/.style={flowbox},
  decision/.style={flowbox, diamond, aspect=3.0, minimum width=3.35cm,
    minimum height=0.95cm},
  io/.style={flowbox, trapezium, trapezium left angle=72,
    trapezium right angle=108, minimum width=1.90cm, minimum height=0.76cm},
]
  \path[use as bounding box] (-1.85,2.20) rectangle (4.40,9.95);

  \node[terminal] (start) at (0,9.55) {开始};
  \node[process, minimum width=2.70cm] (init) at (0,8.25)
    {$k=1,\ s=1$};
  \node[decision] (test) at (0,5.35) {$k<4?$};
  \node[process, minimum width=2.25cm] (update-s) at (3.10,6.45)
    {$s=2s-k$};
  \node[process, minimum width=2.25cm] (update-k) at (3.10,7.45)
    {$k=k+1$};
  \node[io] (output) at (0,3.90) {输出 $s$};
  \node[terminal] (finish) at (0,2.65) {结束};

  % Main sequence; the stem above test is unarrowed and only its lower
  % segment carries the downward arrow.
  \draw[->] (start.south) -- (init.north);
  \coordinate (merge-level) at (0,7.80);
  \draw (init.south) -- (merge-level);
  \draw[->] (merge-level) -- (test.north);
  \draw[->] (output.south) -- (finish.north);

  % Affirmative branch: update s, then update k, then independently return
  % left to the shared stem above the test node.
  \draw[->] (test.east) -- node[above] {是} (3.10,5.35);
  \draw[->] (3.10,5.35) -- (update-s.south);
  \draw[->] (test.south) -- node[right] {否} (output.north);
  \draw[->] (update-s.north) -- (update-k.south);
  \draw[->] (update-k.north) -- (3.10,7.80) -- (merge-level);
\end{tikzpicture}
